\section{Exact certificate and conclusion}
\label{sec:certificate-conclusion}

The analytic proofs above are symbolic.  We include a deterministic
finite certificate to audit the identities and rational exponent
ledger most vulnerable to normalization, boundary, or gcd errors.  It
does not test the logarithmically averaged theorem, estimate primes, or
replace any asymptotic argument.

\subsection{Certificate}

The script
\begin{center}
 \texttt{experiments/tpc23\_certificate.py}
\end{center}
uses only the Python standard library.  It represents
\(-\mu(n)\log n\) exactly as a vector on the formal basis
\(\{\log p:p\text{ prime}\}\), and uses
\texttt{fractions.Fraction} for every exponent calculation.  It checks:
\begin{enumerate}[label=\textup{(\arabic*)}]
 \item complement reflection \eqref{eq:exact-complement} for all
 \(2\le n\le400\) and \(1\le Y\le24\), including nonsquarefree
 targets;
 \item the Liouville cocycle \eqref{eq:exact-Liouville-cocycle} and
 balanced hyperbola split on the same models;
 \item the general CRT compatibility
 \((r,s)\mid h(m_1-m_2)\) and affine determinant
 \eqref{eq:affine-determinant} over five primitive row--shift samples
 and a finite modulus family that includes complements sharing primes
 with \(h\);
 \item a hard-failure variant which omits the factor \((r,s)\) from the
 determinant and is required to disagree;
 \item the dynamic three-cell inequality on an exact rational grid;
 \item both inherited profile collar identities
 \eqref{eq:published-rho} and \eqref{eq:li-rho}.
\end{enumerate}

Independent ordinary and \texttt{python -O} runs perform \(56{,}714\)
exact checks and produce byte-identical output.  The hashes are
\begin{center}
\scriptsize
\begin{tabular}{@{}ll@{}}
\toprule
JSON SHA-256
&\nolinkurl{62e982d5da46d68fcf698f498a2d8fa66cb54ca9422d470f8516bc18f2aa08ab}\\
source SHA-256
&\nolinkurl{51640ec4cbc5c82d792ce40a170ec919ffa2593fade6f92eaffca25d612429be}\\
certificate digest
&\nolinkurl{c5054fcce65ccbdda59dba99c519f5fd9fbe147384f63f6b5d0e17265876be97}\\
\bottomrule
\end{tabular}
\end{center}
The JSON flags for asymptotic evidence, full residual closure,
Hardy--Littlewood, twin primes, and a general parity breakthrough are
all false.

\subsection{Conclusion}

The large-divisor support left after TPC-22 is not an indivisible black
box.  The exact-conductor method is stable under a dynamic residual
cutoff, and its new exponent ledger proves that every strict interior
point of the old hybrid closure region contains a nonzero collar beyond
\(R\).  The enlarged square \(u,v\le S\) includes two mixed pieces and
a both-large piece.  This is a genuine fixed-power analytic advance,
not a numerical observation or a renamed short--short theorem.

At the next boundary, divisor reflection changes the problem.  It
shortens the complementary modulus but transfers the coefficient to
the quotient \(\mu((mj+h)/r)\).  On two rows, the quotient variables
are nonproportional affine forms with exact determinant
\(h(m_1-m_2)/(r,s)\).  This identifies a real two-point parity
cocycle.  Known logarithmic Elliott theory cancels each fixed fiber and
every fixed finite complement family; a diagonal argument even permits
an ineffective slowly growing family.  The absolute-mass theorem shows
why the signs cannot be removed.

What remains is now narrower and more explicit: obtain a rate uniform
over polynomially growing affine coefficients and complement fibers in
the long range, and a summed incidence estimate in the sparse range.
Either advance would move the full residual program forward.  Neither
is supplied here, and neither may be replaced by an assertion of random
M\"obius signs.  Full residual reassembly, a positive prime-pair main
term, and the twin-prime conjecture remain separate open stages.
